% Unit 064 terjemahan Bahasa Indonesia dari chapter5.tex baris 9--140.
% Sumber: Methods of Algebra, Volume 2, commit
% 9a5803ff77dd3257484cb177f851a73770a59dd3 (CC BY 4.0).
% stable-unit-id: o014.aljabr2.chapter5.intro-filtration-grading
% Cakupan: pendahuluan Bab 5 dan Bagian 5.1 lengkap.

% segment-id: o014.li2.chapter5.intro-filtration-grading.h001
\chapter{Barisan Spektral}\label{sec:ss}
\hypertarget{o014.aljabr2.chapter5.intro-filtration-grading}{ }

% segment-id: o014.li2.chapter5.intro-filtration-grading.p001
Barisan spektral merupakan alat aljabar klasik yang terus digunakan dan
dikembangkan, serta berakar kuat dalam topologi. Dalam beberapa buku
ajar atau catatan kuliah, barisan spektral bahkan dipakai sebagai alat dasar
aljabar homologis untuk membangun sebagian sifat fundamental dalam
\CHref{sec:Abel-cat} dan \CHref{sec:cplx}.

% segment-id: o014.li2.chapter5.intro-filtration-grading.p002
Barisan spektral pertama kali diperkenalkan oleh J.\ Leray pada dasawarsa 1940,
lalu memperoleh bentuknya secara bertahap melalui pengembangan oleh J.-L.\
Koszul, H.\ Cartan, dan lainnya. Motivasi utama mereka adalah mengkaji grup
homologi berkas serat. Dalam bahasa aljabar, masalah semula setara dengan
memahami homologi (atau kohomologi) suatu kompleks rantai (atau kompleks)
$X$. Kita ambil kasus kompleks sebagai contoh. Objek $X$ sendiri mungkin sulit
dipahami, tetapi dapat dikaji melalui filtrasi yang sesuai
$\cdots\supset\mathrm{F}^pX\supset\mathrm{F}^{p+1}X\supset\cdots$.
Filtrasi $\mathrm{F}^\bullet X$ menginduksi filtrasi pada setiap $\Hm^n(X)$.
Dari informasi mengenai $\mathrm{F}^\bullet X$ dan subkuosiennya, kita dapat
secara bertahap mendekati setiap subkuosien
$\gr^p\Hm^n(X)$ dari kohomologi terfiltrasi tersebut. Barisan spektral adalah
seni menyusun pendekatan bertahap semacam ini.

% segment-id: o014.li2.chapter5.intro-filtration-grading.p003
Kompleks dapat pula diperumum menjadi objek diferensial $(X,d)$ dalam kategori
abelian (Definisi~\ref{def:differential-object}), sehingga diperoleh definisi
abstrak barisan spektral dalam Definisi~\ref{def:SS-abstract}. Barisan spektral
adalah deret objek diferensial $\mathscr{E}=(E_r,d_r)_r$ sedemikian sehingga
kohomologi setiap ``halaman'' $\Hm(E_r,d_r)$ sama dengan halaman berikutnya
$E_{r+1}$. Jika struktur bergradasi ditambahkan dan diferensial $d$ diizinkan
memiliki derajat, diperoleh konsep barisan spektral bergradasi $E_r^p$ dan
bergradasi ganda $E_r^{p,q}$ (Definisi~\ref{def:graded-spectral-sequence}).
Secara konkret, arah diferensial pada barisan spektral bergradasi ganda adalah
\begin{gather*}
	d_r^{p, q}: E_r^{p, q} \to E_r^{p+r, q-r+1}, \\
	E_r = (E_r^{p, q})_{(p, q) \in \Z^2}, \quad d_r = (d_r^{p, q})_{(p, q) \in \Z^2}.
\end{gather*}
Kompleks terfiltrasi menghasilkan objek diferensial terfiltrasi, lalu secara
alami menghasilkan barisan spektral bergradasi ganda. Inilah bentuk barisan
spektral yang paling sering muncul dalam penerapan. Semua konsep mempunyai
versi yang bersesuaian bagi kompleks rantai dan homologinya.

% segment-id: o014.li2.chapter5.intro-filtration-grading.p004
Strategi bab ini dimulai dari teori pasangan eksak W.\ Massey. Pasangan eksak
merupakan perangkat lain untuk menghasilkan barisan spektral dan cakupannya
lebih luas daripada objek diferensial terfiltrasi. Setelah pengantar filtrasi
dan struktur bergradasi dalam \S\ref{sec:grading-filtration}, kita memberikan
definisi umum barisan spektral dalam \S\ref{sec:spectral-sequence}, termasuk
versi bergradasi dan bergradasi ganda, serta istilah dasar seperti degenerasi,
keterbatasan, dan halaman limit. Catatan~\ref{rem:edge-computing} pada akhir
bagian itu menjelaskan cara membaca barisan eksak dari tepi barisan spektral,
suatu teknik yang penting.

% segment-id: o014.li2.chapter5.intro-filtration-grading.p005
Pasangan eksak didefinisikan dalam \S\ref{sec:exact-couples}. Selanjutnya,
\S\ref{sec:filtered-diff-ss} menjelaskan bagaimana pasangan eksak menghasilkan
barisan spektral bergradasi dari objek diferensial terfiltrasi, dan
\S\ref{sec:filtered-cplx-ss} menjelaskan lebih lanjut bagaimana kompleks
terfiltrasi menghasilkan barisan spektral bergradasi ganda. Di sana juga
dibahas konsep inti yang disebut konvergensi dan
Teorema Konvergensi Klasik~\ref{prop:classical-convergence}. Sebagian
verifikasinya sedikit teknis, tetapi tidak mengandung kesulitan mendasar.
Pembaca dapat terlebih dahulu bekerja dalam kategori abelian konkret seperti
$\cate{Ab}$ untuk menyederhanakan argumen.

% segment-id: o014.li2.chapter5.intro-filtration-grading.p006
Bikompleks menghasilkan kompleks terfiltrasi dengan dua cara, bergantung pada
apakah kompleks total difiltrasi menurut koordinat horizontal atau vertikal.
Barisan spektral yang bersesuaian berturut-turut ditulis
$\mathscr{E}_{\mathrm{I}}$ dan $\mathscr{E}_{\mathrm{II}}$. Perinciannya
merupakan pokok \S\ref{sec:double-cplx-ss}. Kita juga memberikan penerapan
standar, termasuk penurunan bifunktor seimbang
(Contoh~\ref{eg:bifunctor-ss}), barisan spektral funktor hiperturunan
(Contoh~\ref{eg:hyperderived-ss}), dan barisan spektral Grothendieck
(Teorema~\ref{prop:Grothendieck-ss}). Resolusi Cartan--Eilenberg yang telah
diperkenalkan (Teorema~\ref{prop:CE-resolution}) berperan di sini.

% segment-id: o014.li2.chapter5.intro-filtration-grading.p007
Barisan spektral Grothendieck berkaitan dengan penurunan komposisi funktor dan
sangat umum dalam geometri. Sebagai contoh funktor turunan kanan, tinjau
funktor aditif eksak kiri antarkategori abelian
$\mathcal{A}\xrightarrow{F}\mathcal{A}'\xrightarrow{F'}\mathcal{A}''$,
dengan $\mathcal{A}$ dan $\mathcal{A}'$ mempunyai cukup banyak objek
injektif. Jika $F$ memetakan objek injektif ke objek $F'$-asiklik
(Konvensi~\ref{con:F-acyclic}), terdapat barisan spektral bergradasi ganda
kanonik yang konvergen kuat
\[ E_2^{p, q} = (\mathrm{R}^p F') (\mathrm{R}^q F)(X) \Rightarrow \mathrm{R}^{p+q}(F' F)(X), \quad X \in \Obj(\mathcal{A}). \]
Penurunan komposisi funktor juga dapat ditangani pada tingkat kategori turunan;
lihat Teorema~\ref{prop:localization-triangulated-composite}. Versi barisan
spektralnya, selain bersifat elementer secara formal dan mudah dihitung, dalam
beberapa kasus konkret memberikan informasi tambahan seperti barisan eksak
pada derajat rendah.

% segment-id: o014.li2.chapter5.intro-filtration-grading.p008
Bagian terakhir bab ini, \S\ref{sec:multiplicative-SS}, memberikan pengantar
singkat tentang struktur perkalian pada barisan spektral. Aspek ini telah
mendapat perhatian sejak awal sejarah barisan spektral.
Proposisi~\ref{prop:dga-mult-ss} menunjukkan bahwa aljabar diferensial
bergradasi terfiltrasi menghasilkan barisan spektral dengan struktur perkalian
kanonik; pembahasan yang lebih mendalam diserahkan kepada monograf khusus.

% segment-id: o014.li2.chapter5.intro-filtration-grading.p009
Sebagai latihan penerapan, barisan spektral Lyndon--Hochschild--Serre yang
sangat penting dalam teori kohomologi grup dapat diperoleh sebagai penerapan
barisan spektral Grothendieck maupun secara konkret dari kompleks terfiltrasi.
Pendekatan kedua juga memberinya struktur perkalian kanonik. Pokok ini akan
dibahas dalam \S\ref{sec:LHS-SS} pada masa mendatang.

% segment-id: o014.li2.chapter5.intro-filtration-grading.note001
\begin{wenxintishi}
	Bab ini hanya memerlukan pengetahuan mengenai kategori abelian dan kompleks.
Selain beberapa definisi sederhana, bab ini tidak bergantung pada isi
\CHref{sec:triangulated-derived-cat}; kedua bab dapat dibaca secara mandiri.
Karena barisan spektral objek diferensial terfiltrasi dapat ditulis secara
langsung, teori pasangan eksak tidak mutlak diperlukan untuk konstruksi itu
(lihat Proposisi~\ref{prop:filtered-diff-E} dan pembahasan sesudahnya). Pembaca
dapat mempertimbangkan untuk sementara melewati \S\ref{sec:exact-couples}.
Selain itu, di bagian selanjutnya dari buku ini hanya \S\ref{sec:LHS-SS} yang
memerlukan pengetahuan tentang barisan spektral.
\end{wenxintishi}

% segment-id: o014.li2.chapter5.intro-filtration-grading.h002
\section{Filtrasi dan Struktur Bergradasi}\label{sec:grading-filtration}

% segment-id: o014.li2.chapter5.intro-filtration-grading.p010
Kita telah menjumpai objek bergradasi dalam Definisi~\ref{def:graded-obj}.
Objek bergradasi-$\Z$ dalam sembarang kategori $\mathcal{A}$ cukup disebut
\emph{objek bergradasi}; menurut definisi, objek-objek tersebut adalah
$(X^p)_{p\in\Z}$ dalam $\mathcal{A}^{\Z}$. Objek bergradasi-$\Z^2$ cukup
disebut \emph{objek bergradasi ganda}; menurut definisi, objek-objek tersebut
adalah $(X^{p,q})_{(p,q)\in\Z^2}$ dalam
$\mathcal{A}^{\Z\times\Z}$. Morfisme antarobjek bergradasi (atau bergradasi
ganda) seperti biasa ditulis $(f^p)_p$ (atau $(f^{p,q})_{p,q}$), dan
dikomposisikan derajat demi derajat. \index{objek bergradasi}

% segment-id: o014.li2.chapter5.intro-filtration-grading.p011
Jika $\mathcal{A}$ merupakan kategori abelian, maka
$\mathcal{A}^{\Z}$ dan $\mathcal{A}^{\Z\times\Z}$ juga kategori abelian;
kernel, kokernel, dan operasi lain semuanya dihitung derajat demi derajat.

% segment-id: o014.li2.chapter5.intro-filtration-grading.p012
Kategori $\mathcal{A}^{\Z}$ mempunyai autoekuivalensi pergeseran $T$ yang
ditentukan oleh $(TX)^p=X^{p+1}$ dan $(Tf)^p=f^{p+1}$. Lebih umum,
$\mathcal{A}^{\Z^n}$ mempunyai keluarga autoekuivalensi pergeseran
$T_1,\ldots,T_n$ yang bersesuaian dengan pergeseran sepanjang $n$ koordinat.
Pergeseran ini memenuhi komutativitas ketat $T_iT_j=T_jT_i$.

% segment-id: o014.li2.chapter5.intro-filtration-grading.p013
Dalam kerangka bab ini, struktur bergradasi terutama berasal dari filtrasi.

% segment-id: o014.li2.chapter5.intro-filtration-grading.d001
\begin{definisi}\label{def:filtration}
	\index{filtrasi}
	\index{filtrasi!hingga}
	\index[sym1]{FilpX@$\mathrm{F}^p X$, $\mathrm{F}_p X$}
	Tinjau objek $X$ dari kategori aditif $\mathcal{A}$.
	\begin{itemize}
		\item Suatu \emph{filtrasi menurun} $\mathrm{F}^\bullet X$ berarti
		barisan subobjek
		\[ \cdots \supset \mathrm{F}^p X \supset \mathrm{F}^{p+1} X \supset \cdots \quad (p \in \Z). \]
		\item Jika terdapat $a\leq b$ dengan $\mathrm{F}^aX=X$ dan
		$\mathrm{F}^bX=0$, filtrasi tersebut disebut \emph{hingga}.
		\item Filtrasi menaik $\mathrm{F}_\bullet X$ dan sifat hingganya
		didefinisikan secara serupa.
	\end{itemize}
\end{definisi}

% segment-id: o014.li2.chapter5.intro-filtration-grading.p014
Filtrasi menurun dan menaik dapat dibedakan dari posisi indeks $p$. Jika tidak
ada kerancuan, keduanya cukup disebut filtrasi. Dalam kajian kompleks dan
kohomologi biasanya digunakan filtrasi menurun, sedangkan untuk kompleks
rantai dan homologi biasanya digunakan filtrasi menaik. Bab ini terutama
memakai filtrasi menurun.

% segment-id: o014.li2.chapter5.intro-filtration-grading.p015
Semua objek terfiltrasi $(X,\mathrm{F}^\bullet X)$ membentuk kategori
$\mathrm{Fil}^\bullet(\mathcal{A})$. Morfisme
$f:(X,\mathrm{F}^\bullet X)\to(Y,\mathrm{F}^\bullet Y)$ didefinisikan sebagai
morfisme $f:X\to Y$ yang mempertahankan filtrasi, yakni yang untuk setiap $p$
membatasi diri menjadi $\mathrm{F}^pX\to\mathrm{F}^pY$. Kategori
$\mathrm{Fil}_\bullet(\mathcal{A})$ didefinisikan serupa.
\index[sym1]{FilA@$\mathrm{Fil}^\bullet(\mathcal{A})$, $\mathrm{Fil}_\bullet(\mathcal{A})$}

% segment-id: o014.li2.chapter5.intro-filtration-grading.p016
Untuk setiap $X$, filtrasi menurun $\mathrm{F}^\bullet X$ (atau filtrasi
menaik $\mathrm{F}_\bullet X$) menghasilkan objek bergradasi
$(\mathrm{F}^pX)_{p\in\Z}$ (atau $(\mathrm{F}_pX)_{p\in\Z}$). Ini memberikan
funktor $\mathrm{Fil}^\bullet(\mathcal{A})\to\mathcal{A}^{\Z}$ (atau
$\mathrm{Fil}_\bullet(\mathcal{A})\to\mathcal{A}^{\Z}$), yang disebut
konstruksi Rees.

% segment-id: o014.li2.chapter5.intro-filtration-grading.p017
Jika $\mathcal{A}$ kategori abelian, terdapat cara lain membangun objek
bergradasi dari filtrasi.

% segment-id: o014.li2.chapter5.intro-filtration-grading.d002
\begin{definisi}
	\index[sym1]{grX@$\gr^p X$, $\gr_p X$}
	Misalkan $\mathcal{A}$ kategori abelian dan tinjau filtrasi
	$\mathrm{F}^\bullet X$ (atau $\mathrm{F}_\bullet X$) pada objek $X$.
	Definisikan objek bergradasi $\gr X$ melalui
	\[ \gr^p X := \mathrm{F}^p X / \mathrm{F}^{p+1} X \quad (\text{atau} \quad \gr_p X := \mathrm{F}_p X / \mathrm{F}_{p-1} X). \]
	Ini memberikan funktor
	$\gr:\mathrm{Fil}^\bullet(\mathcal{A})\to\mathcal{A}^{\Z}$ (atau
	$\gr:\mathrm{Fil}_\bullet(\mathcal{A})\to\mathcal{A}^{\Z}$).
\end{definisi}

% segment-id: o014.li2.chapter5.intro-filtration-grading.p018
Untuk mengekstrak informasi mengenai $X$ dari $(\gr^pX)_p$, setidaknya kita
memerlukan $\bigcup_p\mathrm{F}^pX=X$ dan
$\bigcap_p\mathrm{F}^pX=0$, dengan notasi dari
Konvensi~\ref{con:increasing-union}. Agar pernyataan pertama bermakna dan
berguna, kita juga menghendaki $\varinjlim_{p\to-\infty}$ eksak dalam
$\mathcal{A}$, atau cukup mensyaratkan adanya $M$ dengan
$\mathrm{F}^MX=X$. Kita mulai dengan beberapa konsep terkait.

% segment-id: o014.li2.chapter5.intro-filtration-grading.d003
\begin{definisi}\label{def:filtration-properties}
	\index{filtrasi!lengkap}
	\index{filtrasi!terpisah}
	\index{filtrasi!menyeluruh}
	Misalkan $\mathcal{A}$ kategori abelian dan $X$ objek dengan filtrasi
	$\mathrm{F}^\bullet X$.
	\begin{itemize}
		\item Jika $\bigcap_p\mathrm{F}^pX=0$, filtrasi tersebut disebut
		\emph{terpisah}.
		\item Jika $\bigcup_p\mathrm{F}^pX=X$ dan salah satu syarat berikut
		berlaku, filtrasi tersebut disebut \emph{menyeluruh}:
		\begin{compactenum}[(a)]
			\item $\mathcal{A}$ mempunyai $\varinjlim$ tersaring terhitung yang
			eksak; lebih tepatnya,
			$\varinjlim:\mathcal{A}^{(\Z_{\geq0},\leq)}\to\mathcal{A}$ ada dan
			eksak\footnote{Jika $\mathcal{A}$ merupakan kategori Grothendieck,
			maka (a) otomatis berlaku.};
			\item terdapat $M\in\Z$ dengan $\mathrm{F}^MX=X$.
		\end{compactenum}
		\item Jika keluarga morfisme kanonik
		$X\twoheadrightarrow X/\mathrm{F}^pX$ menginduksi isomorfisme
		$X\rightiso\varprojlim_pX/\mathrm{F}^pX$, filtrasi tersebut disebut
		\emph{lengkap}.
	\end{itemize}
	Filtrasi menaik $\mathrm{F}_\bullet X$ diperlakukan dengan cara yang sama.
\end{definisi}

% segment-id: o014.li2.chapter5.intro-filtration-grading.p019
Konsep kelengkapan diilhami oleh kasus grup topologis; lihat
\cite[\S 4.10]{Li1}.

% segment-id: o014.li2.chapter5.intro-filtration-grading.p020
Filtrasi hingga otomatis terpisah, menyeluruh, dan lengkap; inilah kasus utama
yang diperlukan dalam bab ini. Perhatikan bahwa jika terdapat $N$ dengan
$\mathrm{F}^NX=0$, filtrasi $\mathrm{F}^\bullet X$ terpisah dan lengkap.
Selain itu, karena kernel dari
$X\to\varprojlim_pX/\mathrm{F}^pX$ adalah
$\bigcap_p\mathrm{F}^pX$, kelengkapan menyiratkan keterpisahan.

% segment-id: o014.li2.chapter5.intro-filtration-grading.p021
Citra filtrasi menyeluruh tetap menyeluruh; hal ini diperoleh dengan
menafsirkan $\bigcup_p$ sebagai $\sum_p$, lalu menggunakan
Lema~\ref{prop:image-sum-preimage-intersection} (ii).
Selain itu, filtrasi menyeluruh memenuhi
\begin{equation}\label{eqn:exhaustion-intersection}\begin{gathered}
	\varinjlim_{p \to -\infty} \mathrm{F}^p X \rightiso \bigcup_p \mathrm{F}^p X = X, \\
	K = K \cap \left( \bigcup_p \mathrm{F}^p X \right) = \bigcup_p \left( K \cap \mathrm{F}^p X \right),
\end{gathered}\end{equation}
dengan $K$ sembarang subobjek dari $X$. Di bawah syarat (a) untuk filtrasi
menyeluruh, pernyataan ini merupakan akibat sederhana dari
Proposisi~\ref{prop:Grothendieck-cat-intersection}; di bawah syarat (b),
pernyataan ini langsung.

% segment-id: o014.li2.chapter5.intro-filtration-grading.pr001
\begin{proposisi}\label{prop:gr-isom}
	Misalkan $\mathcal{A}$ kategori abelian dan
	$f:(X,\mathrm{F}^\bullet X)\to(Y,\mathrm{F}^\bullet Y)$ suatu morfisme dalam
	$\mathrm{Fil}^\bullet(\mathcal{A})$.
	\begin{enumerate}[(i)]
		\item Misalkan $\mathrm{F}^\bullet X$ terpisah dan menyeluruh. Jika
		$\gr(f)$ monomorfisme, maka $f$ monomorfisme.
		\item Misalkan $\mathrm{F}^\bullet X$ lengkap dan menyeluruh, sedangkan
		$\mathrm{F}^\bullet Y$ terpisah dan menyeluruh. Jika $\gr(f)$
		isomorfisme, maka $f$ juga isomorfisme, dan dalam hal ini
		$\mathrm{F}^\bullet Y$ juga lengkap.
	\end{enumerate}
	Pernyataan yang bersesuaian berlaku pula bagi filtrasi menaik.
\end{proposisi}
% segment-id: o014.li2.chapter5.intro-filtration-grading.pf001
\begin{bukti}
	Untuk setiap $p\in\Z$, tuliskan diagram komutatif dengan baris-baris eksak
% segment-id: o014.li2.chapter5.intro-filtration-grading.g001
	\[\begin{tikzcd}
		0 \arrow[r] & \mathrm{F}^{p+1} X \arrow[d, "f"'] \arrow[r] & \mathrm{F}^p X \arrow[r] \arrow[d, "f"] & \gr^p X \arrow[r] \arrow[d, "{\gr^p(f)}"] & 0 \\
		0 \arrow[r] & \mathrm{F}^{p+1} Y \arrow[r] & \mathrm{F}^p Y \arrow[r] & \gr^p Y \arrow[r] & 0
	\end{tikzcd}\]
	Tuliskan $K^p$ dan $C^p$ untuk kernel dan kokernel dari
	$\mathrm{F}^pX\xrightarrow{f}\mathrm{F}^pY$.

	Untuk (i), terapkan Teorema~\ref{prop:snake-lemma} pada diagram di atas.
	Diperoleh $K^{p+1}=K^p$. Secara khusus,
	$K^p\subset\bigcap_{\ell\geq p}\mathrm{F}^\ell X=0$, sehingga pembatasan
	$f$ ke setiap $\mathrm{F}^pX$ monomorfisme. Untuk membuktikan bahwa $f$
	monomorfisme, cukup gunakan
	\[ \Ker(f) \xlongequal{\because\; \text{\eqref{eqn:exhaustion-intersection}}} \bigcup_p \left( \Ker(f) \cap \mathrm{F}^p X \right) = 0. \]

	Untuk (ii), diagram menunjukkan $K^{p+1}=K^p$ dan
	$C^{p+1}\rightiso C^p$. Untuk setiap $\ell>p$, perhatikan diagram komutatif
dengan baris-baris eksak
% segment-id: o014.li2.chapter5.intro-filtration-grading.g002
	\[\begin{tikzcd}
		0 \arrow[r] & \mathrm{F}^{\ell} X \arrow[r] \arrow[d, "f"] & \mathrm{F}^p X \arrow[r] \arrow[d, "f"] & \mathrm{F}^p X / \mathrm{F}^{\ell} X \arrow[r] \arrow[d] & 0 \\
		0 \arrow[r] & \mathrm{F}^{\ell} Y \arrow[r] & \mathrm{F}^p Y \arrow[r] & \mathrm{F}^p Y / \mathrm{F}^{\ell} Y \arrow[r] & 0
	\end{tikzcd}\]
	Dengan menerapkan Teorema~\ref{prop:snake-lemma} sekali lagi, diperoleh
	isomorfisme yang diinduksi oleh $f$
	\[ \mathrm{F}^p X / \mathrm{F}^\ell X \rightiso \mathrm{F}^p Y / \mathrm{F}^\ell Y. \]

	Ambil $\varinjlim_{p:p<\ell}$ pada kedua ruas dan gunakan
	\eqref{eqn:exhaustion-intersection}; diperoleh
	$X/\mathrm{F}^\ell X\rightiso Y/\mathrm{F}^\ell Y$. Selanjutnya ambil
	$\varprojlim_\ell$ untuk mendapatkan diagram komutatif
% segment-id: o014.li2.chapter5.intro-filtration-grading.g003
	\[\begin{tikzcd}
		X \arrow[r, "\sim"] \arrow[d, "f"'] & \varprojlim_\ell X/\mathrm{F}^\ell X \arrow[d, "\sim" sloped] \\
		Y \arrow[hookrightarrow, r] & \varprojlim_\ell Y/\mathrm{F}^\ell Y
	\end{tikzcd}\]
	Maka baik $f$ maupun $Y\to\varprojlim_\ell Y/\mathrm{F}^\ell Y$
	merupakan isomorfisme.
\end{bukti}
